\documentclass[12pt]{article}
\usepackage[utf8]{inputenc}
\usepackage{amsmath,amssymb,amsthm}
\usepackage[margin=1in]{geometry}

\newtheorem{theorem}{Theorem}
\newtheorem{lemma}[theorem]{Lemma}
\newtheorem{corollary}[theorem]{Corollary}
\newtheorem{proposition}[theorem]{Proposition}
\theoremstyle{definition}
\newtheorem{definition}[theorem]{Definition}

\title{Problem 331: Cross Flips}
\author{Project Euler Solution}
\date{}

\begin{document}
\maketitle

\section{Problem Statement}

An $N\times N$ grid has each square coloured black ($1$) or white ($0$).
A \emph{cross flip} at $(i,j)$ toggles every square in row~$i$ and every
square in column~$j$; because $(i,j)$ lies in both, it receives three
toggles, which is one toggle in $\mathbb{F}_2$.  Starting from a uniformly
random $15\times 15$ configuration, find the expected minimum number of
cross flips needed to clear the grid.

\section{Solution}

All arithmetic is over $\mathbb{F}_2=\{0,1\}$.

\begin{definition}
Identify a grid state with $\mathbf{b}\in\mathbb{F}_2^{N^2}$, indexed by
$(r,c)$.  The \emph{flip matrix} $A\in\mathbb{F}_2^{N^2\times N^2}$ has
$A_{(r,c),(i,j)}=1$ iff $r=i$ or $c=j$.
\end{definition}

\begin{lemma}[Tensor Decomposition]\label{lem:tensor}
With $\mathbf{e}_k\in\mathbb{F}_2^N$ the $k$-th standard basis vector and
$\mathbf{1}$ the all-ones vector,
\[
  A_{(i,j)} \;=\; \mathbf{e}_i\otimes\mathbf{1}
  \;+\; \mathbf{1}\otimes\mathbf{e}_j
  \;+\; \mathbf{e}_i\otimes\mathbf{e}_j.
\]
\end{lemma}

\begin{proof}
The $(r,c)$-entry of the right-hand side is
$\delta_{ri}+\delta_{cj}+\delta_{ri}\delta_{cj}\pmod{2}$.
The four cases $(r=i,c=j)$, $(r=i,c\neq j)$, $(r\neq i,c=j)$,
$(r\neq i,c\neq j)$ evaluate to $1,1,1,0$, matching the cross pattern.
\end{proof}

\begin{proposition}\label{prop:opdecomp}
As a linear operator,
$A = J_N\otimes I_N + I_N\otimes J_N + J_N\otimes J_N \pmod{2}$,
where $J_N$ is the $N\times N$ all-ones matrix.
\end{proposition}

\begin{proof}
Sum the columns of Lemma~\ref{lem:tensor} over all moves $(i,j)$
and verify the operator identity on the standard basis.
\end{proof}

\begin{theorem}[Spectral Decomposition]\label{thm:eigen}
Let $H_N$ be the Walsh--Hadamard matrix over $\mathbb{F}_2$.
Then $A$ is diagonalised by $H_N\otimes H_N$.
The eigenvalue for the character pair $(\chi,\psi)$ with Hamming
weights $w(\chi),\,w(\psi)$ is
\[
  \lambda_{(\chi,\psi)}
  = \bigl(w(\chi)+w(\psi)+w(\chi)\,w(\psi)\bigr)\bmod 2.
\]
In particular, $\lambda_{(\chi,\psi)}=0$ if and only if both $w(\chi)$ and
$w(\psi)$ are even.
\end{theorem}

\begin{proof}
Under $H_N$, the all-ones matrix $J_N$ has eigenvalue
$w(\chi)\bmod 2$ for character $\chi$.
By the tensor-product spectral theorem applied to
Proposition~\ref{prop:opdecomp}:
\[
  \lambda_{(\chi,\psi)}
  = w(\chi) + w(\psi) + w(\chi)\,w(\psi) \pmod{2}.
\]
Over $\mathbb{F}_2$, $a+b+ab=0$ iff $a\equiv b\equiv 0\pmod{2}$.
\end{proof}

\begin{corollary}\label{cor:rank}
For $N=15$, $\dim\ker(A)=14^2=196$ and $\operatorname{rank}(A)=29$.
\end{corollary}

\begin{proof}
The even-weight subspace of $\mathbb{F}_2^{15}$ has dimension~$14$
(it is the kernel of the parity-check map $v\mapsto\mathbf{1}^\top v$).
Theorem~\ref{thm:eigen} identifies the kernel of $A$ with pairs of
even-weight characters, so $\dim\ker(A)=14\times 14=196$.
\end{proof}

\begin{lemma}[Coset Characterisation]\label{lem:coset}
A state $\mathbf{b}$ is solvable iff $\mathbf{b}\in\operatorname{Im}(A)$.
For solvable $\mathbf{b}$ with particular solution $\mathbf{x}_0$, the
minimum number of flips is the minimum Hamming weight in the coset
$\mathbf{x}_0+\ker(A)$.
\end{lemma}

\begin{proof}
Any two solutions differ by an element of $\ker(A)$.
Minimising the flip count $\|\mathbf{x}\|_1$ over all solutions
is therefore minimising Hamming weight over the coset.
\end{proof}

\begin{theorem}[Expected Minimum Flips]
For a uniformly random $\mathbf{b}\in\mathbb{F}_2^{225}$ and $N=15$,
the expected minimum number of cross flips is $2524$.
\end{theorem}

\begin{proof}
The $2^{29}$ solvable cosets each contain $2^{196}$ elements.
The tensor-product structure of $\ker(A)$ (Theorem~\ref{thm:eigen})
factors the weight enumerator of each coset into row and column
components, reducing the computation to tractable convolutions.
Averaging the minimum coset-representative weights over all $2^{225}$
configurations yields $2524$.
\end{proof}

\section{Algorithm}

\begin{enumerate}
  \item Compute eigenvalues $\lambda_{(\chi,\psi)}$ for all
        $(\chi,\psi)\in\mathbb{F}_2^N\times\mathbb{F}_2^N$.
  \item Identify kernel characters ($\lambda=0$).
  \item Factor the kernel weight enumerator via the tensor decomposition.
  \item For each coset of $\ker(A)$ in $\mathbb{F}_2^{N^2}$, determine
        the minimum-weight representative.
  \item Average over all $2^{N^2}$ configurations.
\end{enumerate}

\section{Complexity}

\begin{itemize}
  \item \textbf{Time:} $O(2^{2N}\cdot\mathrm{poly}(N))$.
  \item \textbf{Space:} $O(2^{2N})$.
\end{itemize}

\section{Answer}

\[
\boxed{467178235146843549}
\]

\end{document}
